\Legend
对固定座位，每人恰好把消息传给一个后继，因此形成一个有向函数图。若某个座位入度为零，只有它被选为初始接收者才可能收到消息。所以要从一个起点走遍全部座位，入度为零的点最多只能有一个。

令 $b_i=1$ 表示该座位放数字 $3$，否则为 $0$。能传给座位 $j$ 的只可能是 $j-2$ 上的数字 $2$ 或 $j-3$ 上的数字 $3$，故其入度为 $1-b_{j-2}+b_{j-3}$（下标模 $n$）。入度零恰好对应循环序列 $b$ 中的一次 $1\to0$ 切换。若两种卡片都有，则 $b$ 在环上必须只有一次这样的切换，也就是所有 $3$ 连续成一块、所有 $2$ 连续成另一块。任意这种排列只差整体旋转，因此只需测试一种：先放全部 $3$，再放全部 $2$。

构造后计算入度。混合情形从唯一入度零的点沿后继走，若走到重复点前访问了 $n$ 个点即输出 \tt{Yes}。纯一种卡片时没有入度零的点，从座位 $0$ 出发即可；这也等价于检查 $\gcd(n,2)=1$ 或 $\gcd(n,3)=1$。

参考解换一种验证方式：同一排列的无向底图若连通，则该函数图只有一个弱连通分量，结合上述入度结构，存在覆盖全部点的传递过程。

总时间 $O(\sum n)$，单组额外空间 $O(n)$。
